\DocumentMetadata{
  pdfversion=2.0,
  pdfstandard=ua-2,
  lang=en-US,
  tagging=on,
  tagging-setup={math/setup=mathml-SE,tabsorder=structure}
}
\documentclass[11pt,letterpaper]{article}
\usepackage[margin=0.68in,includefoot]{geometry}
\usepackage{newcomputermodern}
\usepackage{amsmath,amscd,mathtools}
\providecommand{\square}{\mdlgwhtsquare}
\usepackage[hidelinks]{hyperref}
\hypersetup{
  pdftitle={Topology First-Year Exam, August 2023, Part 2},
  pdfauthor={University of Florida Department of Mathematics},
  pdfsubject={Graduate mathematics examination},
  pdfdisplaydoctitle=true,
  bookmarksopen=true
}
\setcounter{secnumdepth}{0}
\setlength{\parindent}{0pt}
\setlength{\parskip}{0.28em}
\setlength{\emergencystretch}{3em}
\setlength{\leftmargini}{2.15em}
\setlength{\leftmarginii}{2.65em}
\setlength{\leftmarginiii}{3em}
\renewcommand{\labelenumi}{\textbf{\theenumi.}}
\pagestyle{empty}
\raggedbottom
\allowdisplaybreaks
\newenvironment{examproblems}[1]{%
  \begin{enumerate}%
  \setcounter{enumi}{#1}%
  \setlength{\topsep}{0.35em}%
  \setlength{\partopsep}{0pt}%
  \setlength{\itemsep}{0.48em}%
  \setlength{\parsep}{0.18em}%
}{\end{enumerate}}
\newenvironment{examsubparts}{%
  \begin{enumerate}%
  \setlength{\topsep}{0.18em}%
  \setlength{\partopsep}{0pt}%
  \setlength{\itemsep}{0.16em}%
  \setlength{\parsep}{0.1em}%
}{\end{enumerate}}
\newenvironment{examinstructions}{%
  \par\begingroup\itshape
}{%
  \par\endgroup\vspace{0.18em}
}
\makeatletter
\renewcommand\section{\@startsection{section}{1}{0pt}%
  {-1.25ex plus -.25ex minus -.1ex}%
  {0.55ex plus .1ex}%
  {\normalfont\large\bfseries}}
\renewcommand{\@maketitle}{%
  \newpage\null\vskip -1.2em%
  {\footnotesize\bfseries
    \href{https://www.ufl.edu/}{University of Florida}, \href{https://math.ufl.edu/}{Department of Mathematics}\par}%
  \vskip 0.18em%
  \tagstructbegin{tag=Title}%
    {\LARGE\bfseries\href{https://gma.math.ufl.edu/exams/topology-first-year/}{Topology First-Year Exam}, August 2023, Part 2\par}%
  \tagstructend%
  \vskip 0.35em%
  \tagmcbegin{artifact}%
    \hrule height 1pt%
  \tagmcend%
  \vskip 0.55em%
}
\makeatother
\title{Topology First-Year Exam, August 2023, Part 2}
\author{}
\date{}
\begin{document}
\maketitle
\thispagestyle{empty}
\begin{examinstructions}
Work the following problems and show all work. Support all statements to the best of your ability.
\end{examinstructions}
\begin{examproblems}{0}
\item
Show that \(B^2\) is not homeomorphic to \(B^3\), where \(B^n\) denotes the closed unit ball in \(\mathbb{R}^n\).
\item
Show that every continuous map \(f:S^2\to S^1\) of the 2-sphere to the circle \(S^1\) is null-homotopic.
\item
Let 
\[
X=([-1,1]\times\{0\}\times\{0\})\cup(\{0\}\times[-1,1]\times\{0\})\cup(\{0\}\times\{0\}\times[-1,1])\subset\mathbb{R}^3.
\]
 Show that every continuous map \(f:X\to X\) has a fixed point.
\item
Derive the Brouwer Fixed Point Theorem for \(B^2\) from the Borsuk–Ulam Theorem for \(S^2\).
\item
Let~\(X\), \(Y\), and~\(Z\) denote the following open rays in \(\mathbb{R}^3\): 
\[
X=(0,\infty)\times\{0\}\times\{0\},\quad Y=\{0\}\times(0,\infty)\times\{0\},\quad Z=\{0\}\times\{0\}\times(0,\infty).
\]
 Compute \(\pi_1(A)\) where
\begin{examsubparts}
\item[\textnormal{(a)}]
\(A=\mathbb{R}^3\setminus X\)?
\item[\textnormal{(b)}]
\(A=\mathbb{R}^3\setminus\bar X\)? Here \(\bar X\) denotes the closure of~\(X\) in \(\mathbb{R}^3\).
\item[\textnormal{(c)}]
\(A=\mathbb{R}^3\setminus(X\cup Y)\)?
\item[\textnormal{(d)}]
\(A=\mathbb{R}^3\setminus(\bar X\cup Y)\)?
\item[\textnormal{(e)}]
\(A=\mathbb{R}^3\setminus(X\cup Y\cup Z)\)?
\item[\textnormal{(f)}]
\(A=\mathbb{R}^3\setminus(\bar X\cup Y\cup Z)\)?
\end{examsubparts}
\end{examproblems}
\begin{examinstructions}
Answer the following with complete definitions or proofs.
\end{examinstructions}
\begin{examproblems}{5}
\item
Describe a space~\(X\) with \(\pi_1(X)\cong\mathbb{Z}_3\).
\item
Show that the projective plane \(\mathbb{R}P^2\) is not homeomorphic to the surface of genus 2, \(M_2=T\#T\).
\item
Let \(X=[0,1]^{(0,1)}\) be given the product topology. Is the space~\(X\)
\begin{examsubparts}
\item[\textnormal{(a)}]
compact?
\item[\textnormal{(b)}]
metrizable?
\item[\textnormal{(c)}]
separable?
\end{examsubparts}
\item
State the Seifert–van Kampen Theorem. Can it be used to compute \(\pi_1(S^n)\)?
\item
State the Tychonoff Theorem. Prove that every compact Hausdorff space~\(X\) can be embedded into a Tychonoff cube \([0,1]^A\).
\end{examproblems}
\end{document}
